\documentclass[a4paper, 12pt]{article}

% Русский язык
\usepackage[T2A]{fontenc}
\usepackage[utf8]{inputenc}
\usepackage[english,russian]{babel}

% Картинки
\usepackage{graphicx}
\graphicspath{{images/}}
\DeclareGraphicsExtensions{.pdf,.png,.jpg}
\usepackage{wrapfig}

% Математика
\usepackage{amsmath,amsfonts,amssymb,amsthm,mathtools} 

% Параметры страницы
\usepackage[left=2cm,right=2cm,top=2cm,bottom=3cm,bindingoffset=0cm]{geometry}

\usepackage{enumitem,xparse}
\usepackage{caption}
\usepackage{subcaption}
\usepackage{float}

\usepackage{hyperref}

\begin{document}

\newgeometry{left=2cm,right=2cm,top=2cm,bottom=1cm,bindingoffset=0cm}

\begin{titlepage}
    \begin{center}

        \vspace*{5cm}

        {\Huge Московский Физико-Технический Институт}\\[2cm]

        {\LARGE Отчёт по эксперименту}\\[0.25cm]

        \noindent\rule{\textwidth}{1pt}
        \vspace*{-0.25cm}

        {\huge \textbf{Взаимодействие тио-сульфата натрия с серной кислотой}}

        \noindent\rule{\textwidth}{1pt}

        \vfill

    \end{center}

    \begin{flushright}
        \begin{minipage}{0.45\textwidth}
            {\Large
                Выполнил:\\
                Студент 1 курса ФАКТ\\
                Группа Б03-504\\
                Подмосковнов Лев
            }
        \end{minipage}
    \end{flushright}

    \vfill

    \begin{center}
        Долгопрудный\\2026
    \end{center}

\end{titlepage}
\restoregeometry

\setcounter{page}{2}

\subsection*{Цель работы}

\begin{itemize}[label=--]
    \item Определить $k$ для $27^\circ C$ и $37^\circ C$
    \item Определить энергию активации $E_a$
    \item Определить предэкспаненту $A$
    \item Определить порядок и молекулярность процесса
\end{itemize}

\subsection*{Реактивы и оборудование}

\begin{itemize}[label=--]
    \item Калориметр
    \item 10 пробирок
    \item Бюретка
    \item $Na_2 S_2 O_3$, $H_2 S O_4$, $H_2 O_{\text{дист}}$
\end{itemize}

\subsection*{Уравнения реакций}
\[
    \text{Na}_2 \text{S}_2 \text{O}_3 + \text{H}_2 \text{S} \text{O}_4  \rightarrow \text{H}_2 \text{S}_2 \text{O}_3  + \text{Na}_2 \text{S} \text{O}_4 \text{ быстро}
\]
\[
    \text{H}_2 \text{S}_2 \text{O}_3 = \text{H}_2 \text{S} \text{O}_3 + \text{S} \downarrow \text{ медленная}
\]
\[
    \text{H}_2 \text{S} \text{O}_3 = \text{H}_2 \text{O} + \text{S} \text{O}_2 \text{ быстро}
\]

\begin{table}[h!]
    \centering
    \begin{tabular}{|c|c|c|c|} % 4 колонки, выравнивание по центру, вертикальные линии
        \hline
        N & $\tau$, с & $T$, $^\circ$C & $c$, $\text{моль}$/$\text{л}$ \\
        \hline
        1 & 180       & 27             & 0.01                          \\
        \hline
        2 & 100       & 27             & 0.02                          \\
        \hline
        3 & 60        & 27             & 0.03                          \\
        \hline
        4 & 45        & 27             & 0.04                          \\
        \hline
        5 & 35        & 27             & 0.05                          \\
        \hline
        1 & 120       & 37             & 0.01                          \\
        \hline
        2 & 70        & 37             & 0.02                          \\
        \hline
        3 & 35        & 37             & 0.03                          \\
        \hline
        4 & 25        & 37             & 0.04                          \\
        \hline
        5 & 20        & 37             & 0.05                          \\
        \hline
    \end{tabular}
    \caption{Пример таблицы на 4 столбца}
    \label{tab:my_table}
\end{table}

\clearpage

\begin{figure}[h!]
    \centering
    \includegraphics[width=0.8\textwidth]{1.png}
    \caption{График зависимости v от c}
\end{figure}

\subsection*{Определить $k$}

\hspace{0.5cm} $k_1 = 0.46 \pm 0.07$ $\frac{\text{л}\cdot \text{моль}}{\text{с}}$

$k_2 = 1.09 \pm 0.08$ $\frac{\text{л}\cdot \text{моль}}{\text{с}}$

$\frac{k_2}{k_1} = 2.37$

Выполнено правило Вант-Гоффа

\subsection*{Определить энергию активации $E_a$}

\hspace{0.5cm} Уравнение Аррениуса: $k=A\cdot e^{-\frac{E_a}{RT}}$

$E_a = \frac{RT_1T_2\ln \frac{k_2}{k_1}}{T_2-T_1}=62.4 \pm 4.2$ кДж

\subsection*{Определить предэкспаненту A}

\hspace{0.5cm} $A = (1.33 \pm 0.6) \cdot 10^{-11}$ $\frac{\text{л}\cdot \text{моль}}{\text{с}}$

\subsection*{Определить порядок и молекулярность процесса}

\hspace{0.5cm} $V(c) $ - линеен

$n = 1$ при $\text{Na}_2 \text{S}_2 \text{O}_3$

Если зафиксировать $\text{Na}_2 \text{S}_2 \text{O}_3$ $n = 1$ при $\text{H}_2 \text{S} \text{O}_4$

$n_{\text{общ}} = 2$

молекулярность лимитированной стадии 1 = молекулярность всей реакции

\section*{Вывод}

Былап исследована зависимосьт скорости реакции от концентрации вещества, проверено правило Вант-Гоффа и посчитаны порядок и молекулярность реакции

\end{document}